% Noether Paper 7 complete zh-Hans-CN translated producer draft.
% Translation source: exact current-German interval, SHA-256
% F6C923B79406542E3DE64298DCD38887FF9A52141C71B8FF2BEBE6D14625FAEA.
% Inherited Simplified-Chinese interval is a drafting witness only, SHA-256
% BB4686153D7241CD0F8A74164B6486C31C3BF731722334CC25B5E81AA8884AF8.
% Status: translated producer draft; independent check pending.
% Compilation is mechanical production only. No source, semantic, formula,
% terminology, visual, regional, publication, archive, or certification check is claimed.
\documentclass[11pt]{article}
\usepackage[a4paper,margin=2.4cm]{geometry}
\usepackage{fontspec}
\usepackage{xeCJK}
\usepackage{amsmath,amssymb,mathtools,array}
\setmainfont{Noto Serif}
\setCJKmainfont{Microsoft YaHei}
\setCJKsansfont{Microsoft YaHei}
\XeTeXlinebreaklocale "zh"
\XeTeXlinebreakskip=0pt plus 1pt
\setlength{\parindent}{2em}
\setlength{\parskip}{0.35em}
\setlength{\emergencystretch}{8em}
\allowdisplaybreaks
\raggedbottom
\providecommand{\Hgrp}{\mathfrak{H}}
\providecommand{\NoetherSrcNote}[2]{\textsuperscript{#1}\begingroup\renewcommand{\thefootnote}{}\footnotetext{\textsuperscript{#1}~#2}\endgroup}
\providecommand{\srcnumdisplay}[3][3.4em]{%
  \par\smallskip\noindent
  \begin{minipage}[t]{#1}\vspace*{0pt}#2\end{minipage}%
  \begin{minipage}[t]{\dimexpr\linewidth-#1\relax}\vspace*{0pt}\centering
  \(\displaystyle #3\)
  \end{minipage}\par\smallskip}
\begin{document}
\setcounter{footnote}{0}

\section*{7. 有限群不变量的有限性定理}
\begin{center}
Math. Ann. 77 (1916), S. 89--92
\end{center}

\begin{center}
{\Large 有限群不变量的有限性定理。}\par
\vspace{1.2em}
作者\par
\vspace{0.8em}
{\scshape Emmy Noether}，埃朗根。
\end{center}

以下将给出有限群不变量有限性的一个完全初等的证明 --- 它仅以对称函数理论为基础 ---，同时还实际给出完整不变量系；而通常的证明依赖 Hilbert 关于模基的定理（Ann. 36），因而只是存在性证明。\NoetherSrcNote{*)}{参见，例如 Weber，Lehrbuch der Algebra（第 2 版），第 2 卷，\S\ 57。}

设有限群 $\Hgrp$ 由 $h$ 个线性变换（其行列式非零）$A_1\ldots A_h$ 组成，其中 $A_k$ 表示线性变换
\[
  x_1^{(k)}=\sum_{\nu=1}^{n}a_{1\nu}^{(k)}x_\nu,\ldots,
  x_n^{(k)}=\sum_{\nu=1}^{n}a_{n\nu}^{(k)}x_\nu
\]
或简记为：$(x^{(k)})=A_k(x)$。因此，群 $\Hgrp$ 将以 $x_1\ldots x_n$ 为元素的变量组 $(x)$ 变换为以 $x_1^{(k)}\ldots x_n^{(k)}$ 为元素的变量组 $(x^{(k)})$。由于 $A_1\ldots A_h$ 中必定包含恒等变换，变量组 $(x^{(k)})$ 中也包含变量组 $(x)$。所谓该群的一个整有理（绝对）不变量，是指 $x_1\ldots x_n$ 的这样一个整有理函数：施以 $A_1\ldots A_h$ 后，它恒等地保持不变；因此，对于这样一个不变量 $f(x)$，有：
\srcnumdisplay{(1)}{%
  f(x)=f(x^{(1)})=\cdots=f(x^{(h)})=
  \frac1h\cdot\sum_{k=1}^{h}f(x^{(k)}).}

\noindent\textbf{1.}\enspace 公式 (1) 表明，$f(x)$ 是变量组 $(x^{(1)})\ldots(x^{(h)})$ 的整有理\emph{对称}函数 --- 而且，由于每一项 $f(x^{(k)})$ 只含\emph{一个}变量组 $(x^{(k)})$，这里所涉及的是最简单的、通常称为\emph{单式}的情形。根据关于变量组的对称函数的已知定理\NoetherSrcNote{*)}{参见第 2 节的注。}，$f(x)$ 因而可由这些变量组的初等对称函数整有理地表示，即由“Galois 预解式”的系数 $G_{\alpha\alpha_1\ldots\alpha_n}(x)$ 表示：
\[
\begin{aligned}
  \Phi(z,u)
  &=\prod_{k=1}^{h}\bigl(z+u_1x_1^{(k)}+\cdots+u_nx_n^{(k)}\bigr) \\
  &=z^h+
    \sum G_{\alpha\alpha_1\ldots\alpha_n}(x)
    z^{\alpha}u_1^{\alpha_1}\cdots u_n^{\alpha_n},
\end{aligned}
\qquad
\left(\begin{array}{c}
\alpha+\alpha_1+\cdots+\alpha_n=h\\
\alpha\ne h
\end{array}\right).
\]
其中 $G_{\alpha\alpha_1\ldots\alpha_n}(x)$ 是关于 $x$ 的次数为 $\alpha_1+\cdots+\alpha_n$ 的不变量。由此证明：

\emph{Galois 预解式的系数 $G_{\alpha\alpha_1\ldots\alpha_n}(x)$ 构成该群的一个完整不变量系，使得每个不变量都可由这有限多个不变量整有理地表示。}
\noindent\textbf{2.}\enspace 也可以从 (1) 出发，作如下更为初等的考察，\NoetherSrcNote{**)}{这归结为在上述\emph{单式}情形下证明关于变量组的对称函数的定理。}这一考察同时导出第二个完整不变量系。令
\[
  f(x)=a+b x_1^{\mu_1}\cdots x_n^{\mu_n}
       +\cdots+c x_1^{\nu_1}\cdots x_n^{\nu_n},
\]
其中 $a,b,\ldots,c$ 表示常数，则由 (1) 有：
\[
\begin{aligned}
 h\cdot f(x)
 &=h\cdot a+b\cdot\sum_{k=1}^{h}(x_1^{(k)})^{\mu_1}\cdots(x_n^{(k)})^{\mu_n}
     +\cdots
     +c\cdot\sum_{k=1}^{h}(x_1^{(k)})^{\nu_1}\cdots(x_n^{(k)})^{\nu_n} \\
 &=h\cdot a+b\cdot J_{\mu_1\ldots\mu_n}+\cdots+c\cdot J_{\nu_1\ldots\nu_n}.
\end{aligned}
\]
因此，每个不变量都可以由这些特殊不变量
\[
  J_{\mu_1\ldots\mu_n}
  =\sum_{k=1}^{h}(x_1^{(k)})^{\mu_1}\cdots(x_n^{(k)})^{\mu_n}
\]
作整线性组合，因而只需对这些特殊不变量加以证明。除一个数因子外，$J_{\mu_1\ldots\mu_n}$ 正是表达式
\[
  S_\mu=\sum_{k=1}^{h}
  \bigl(u_1x_1^{(k)}+\cdots+u_nx_n^{(k)}\bigr)^{\mu},
\]
中 $u_1^{\mu_1}\cdots u_n^{\mu_n}$ 的系数；这里令 $\mu=\mu_1+\cdots+\mu_n$，而该表达式表示 $h$ 个线性形式
\[
  \xi_1=u_1x_1^{(1)}+\cdots+u_nx_n^{(1)},\ldots,
  \xi_h=u_1x_1^{(h)}+\cdots+u_nx_n^{(h)}
\]
的第 $\mu$ 次幂和。众所周知，无穷多个幂和 $S_\mu$ 都是
\[
  S_1,S_2,\ldots,S_h,
\]
的整有理组合，而后者的系数由不变量
\[
  J_{\mu_1\ldots\mu_n}\qquad(\mu_1+\cdots+\mu_n\leqq h)
\]
给出；由此得到第二个完整不变量系：

\emph{该群的一个完整不变量系由所有不变量 $J_{\mu_1\ldots\mu_n}$ 给出，其中 $\mu_1+\cdots+\mu_n\leqq h$，而 $h$ 表示群的阶。}

它同 1. 的联系可由如下观察得出：可以用 $\xi_1,\ldots,\xi_h$ 的初等对称函数代替幂和 $S_1,\ldots,S_h$，由此即得到那里给出的由 $G_{\alpha\alpha_1\ldots\alpha_n}(x)$ 构成的完整不变量系。两个结果都表明，\emph{所有不变量都可由那些关于诸 $x$ 的次数不超过群阶 $h$ 的不变量作整有理表示。}

\noindent\textbf{3.}\enspace 由刚刚证明的结果，可以得到关于有理表示的推论。每一个\emph{有理}绝对不变量都能表示为两个不必互素的整有理绝对不变量之商；这一点只要把分子和分母同乘以分母相对于该群的各个共轭，就能立即看出。因此，\emph{每一个有理绝对不变量都可以由 Galois 预解式 $\Phi(z,u)$ 的系数 $G_{\alpha\alpha_1\ldots\alpha_n}(x)$ 有理表示。}这个定理也可以不借助有限性定理而用多种方式轻易证明；Weber II，\S\ 58 中已经表述了这个定理。\NoetherSrcNote{*)}{不过，那里给出的证明并不成立；它只证明了公式 (7) 中的函数 $\Psi(t)$ 以不变量为系数，却没有证明这些不变量可以由 $\Phi(t)$ 的系数有理表示。若用 $\xi_k$ 表示 $x_i^{(k)}$ 时，以一个已知的微分法代替 Lagrange 插值公式，就可以避开这个缺口。这个方法给出下列关系式，它是将恒等式 $\Phi(-\xi_k,u)=0$ 对所有 $u$ 求导所得：
\[
 \left[\frac{\partial\Phi}{\partial u_i}
      -x_i^{(k)}\cdot\frac{\partial\Phi}{\partial z}\right]_{z=-\xi_k}=0.
\]
因此，对于每一个有理函数 $\omega(x)$，Weber 的公式 (7) 和 (8) 应当由下式代替：
\[
 \omega(x_1^{(k)}\cdots x_n^{(k)})=
 \omega\!\left(
   \frac{\partial\Phi/\partial u_1}{\partial\Phi/\partial z}\cdots
   \frac{\partial\Phi/\partial u_n}{\partial\Phi/\partial z}
 \right)_{z=-\xi_k},
\]
这个式子只含有 $\Phi(z,u)$ 的系数；当 $\omega(x)$ 为不变量时，将它对所有 $k$ 求和，就得到所要求的表示。}

\newpage
\noindent\textbf{4.}\enspace 最后需要指出，这里推导出的结果已隐含在我的论文“有理函数的域与系统”\NoetherSrcNote{*)}{Math. Ann. 76，第 161 页（1915 年）。}中；具体而言，1. 中给出的完整系统包含在定理 VI 和 VII 中，而独立于这个\emph{有限性定理}的有理表示证明包含在定理 III 中。\NoetherSrcNote{**)}{对于\emph{相对}不变量，定理 VIII 和 IX 中包含了一个有限性证明；这个证明所依据的基础也不同于通常的证明，但它同样只是存在性证明。这是因为相对不变量不构成一个域。}不过，与一般理论相比，在这里所讨论的特殊情形中，证明过程和结果都显著简化了。通过与 E. Fischer 先生的交谈，我想到把一般研究应用于有限群的不变量；2. 中给出的证明在内容上也出自他——在\emph{一般}情形中，那里出现的完整系统不能有理定义——3. 的注释同样出自他。

\textbf{埃朗根，1915 年 5 月。}

\clearpage

\end{document}
